% Requires \usepackage{tikz} and \usetikzlibrary{arrows.meta} in the preamble.
\begin{figure}[t]
\centering
\begin{tikzpicture}[
   kept/.style = {circle, draw, thick, fill=black!8, minimum size=7mm, inner sep=0pt, font=\small},
   drop/.style = {circle, draw=black!40, densely dashed, text=black!45, minimum size=7mm, inner sep=0pt, font=\small},
   dep/.style  = {-{Stealth[length=4.5pt]}, draw=black!65},
   ann/.style  = {font=\small, align=center}
]

% (a) uncontracted: chain of dropped vertices between inputs and outputs
\node[kept] (o1) at (0,4.8) {$o_1$};
\node[kept] (o2) at (1.5,4.8) {$o_2$};
\node[kept] (o3) at (3,4.8) {$o_3$};
\node[drop] (w)  at (1.5,3.6) {$w$};
\node[drop] (v)  at (1.5,2.4) {$v$};
\node[drop] (u)  at (1.5,1.2) {$u$};
\node[kept] (i1) at (0,0) {$i_1$};
\node[kept] (i2) at (1.5,0) {$i_2$};
\node[kept] (i3) at (3,0) {$i_3$};

\foreach \o in {o1,o2,o3} \draw[dep] (\o) -- (w);
\draw[dep] (w) -- (v);
\draw[dep] (v) -- (u);
\foreach \i in {i1,i2,i3} \draw[dep] (u) -- (\i);

\node[ann] at (1.5,-1.0) {(a) uncontracted\\$|E| = 8$, \ $o \rightsquigarrow i$ path $= 4$};

% (b) contracted: u, v, w spliced out, each o wired to the frontier {i1,i2,i3}
\begin{scope}[xshift=7cm]
   \node[kept] (p1) at (0,4.8) {$o_1$};
   \node[kept] (p2) at (1.5,4.8) {$o_2$};
   \node[kept] (p3) at (3,4.8) {$o_3$};
   \node[kept] (q1) at (0,0) {$i_1$};
   \node[kept] (q2) at (1.5,0) {$i_2$};
   \node[kept] (q3) at (3,0) {$i_3$};

   \foreach \o in {p1,p2,p3} \foreach \i in {q1,q2,q3} \draw[dep] (\o) -- (\i);

   \node[ann] at (1.5,-1.0) {(b) contracted\\$|E| = 9$, \ $o \rightsquigarrow i$ path $= 1$};
\end{scope}

\draw[-{Stealth[length=6pt]}, thick] (4.3,2.4) -- node[above, font=\small] {contract} (5.7,2.4);

\end{tikzpicture}
\caption{Contraction shortens input--output paths but can multiply edges. Edges point from a
         vertex to the vertices it depends on, so backward slicing reads downwards. Solid
         vertices are kept (inputs, outputs and \texttt{@doc}-annotated values); dashed
         vertices are dropped. Dropping $\alpha$ splices it out by rewiring $\alpha$'s consumers
         to $\alpha$'s frontier --- its nearest kept ancestors --- so in (b) every $o_j$ points
         straight at $\{i_1,i_2,i_3\}$. The chain $w,v,u$ collapses to a single hop, but its
         $3 + 2 + 3 = 8$ edges become $3 \times 3 = 9$. In general a chain of $L$ dropped
         vertices between $p$ inputs and $s$ outputs replaces $p + s + L - 1$ edges with
         $p \cdot s$, an increase whenever $p \cdot s > p + s + L - 1$.}
\label{fig:contraction}
\end{figure}
